% The front matter as data: title, authors, their e-mail, ORCID and contributions, and the one affiliation.
% No journal's syntax appears here, because every class asks for the same facts in a different shape: one wants
% a single name per author with the address as key=value pairs, another wants the given and family names split
% into separate commands and the address broken into fields of its own, and the reading copy wants one
% \and-separated line. So this file states each fact once and a template renders it, exactly as the abstract and
% the keywords already work. Supporting one more journal means writing the renderers below in the new
% template, and changing nothing here.
%
% A template must define, before it uses \projauthorlist:
%   \projauthor{initials}{given}{family}{email}{orcid}{CRediT roles}      one non-corresponding author
%   \projauthorcorr{...}                                                 same six, for the corresponding author
% ORCID may be empty, so a renderer that prints it guards the field with \projifhasorcid.

\newcommand{\projtitle}{Leakage-Free Benchmarking of Classical Machine Learning and Graph Neural Networks for Low-Data TgDHFR Bioactivity Prediction}
\newcommand{\projshorttitle}{Leakage-Free Benchmarking of Classical ML and Graph Neural Networks}
\newcommand{\projrunningauthors}{Vinajera-Zamora et al., 2026}

\newcommand{\projaffildiv}{Instituto Universitario de Investigaci\'on y Desarrollo Tecnol\'ogico}
\newcommand{\projaffilorg}{Universidad Tecnol\'ogica Metropolitana}
\newcommand{\projaffilstreet}{Ignacio Valdivieso 2409}
\newcommand{\projaffilcity}{Santiago}
\newcommand{\projaffilstate}{Regi\'on Metropolitana}
\newcommand{\projaffilpostcode}{97279}
\newcommand{\projaffilcountry}{Chile}

% Not every author has an ORCID on file, and a renderer that prints the field unconditionally would emit an empty
% link for the ones who do not. This runs its second argument only when the first is a non-empty ORCID.
\newcommand{\projifhasorcid}[2]{\def\projorcidarg{#1}\ifx\projorcidarg\empty\else#2\fi}

% The authors in the order they appear on the paper. The corresponding author is marked by which macro the line
% calls, not by a flag, so no layout has to decide what "corresponding" looks like in its own notation.
% The roles are the CRediT taxonomy's own terms: they go into the author block of a class that has one, and into
% the contributions declaration, which is rendered from this same list.
\newcommand{\projauthorlist}{%
	\projauthorcorr{A.V.-Z.}{Andrey}{Vinajera-Zamora}{avinajera@utem.cl}{0000-0002-4922-457X}%
		{Conceptualization, Methodology, Formal analysis, Software, Validation, Investigation,
		 Writing -- original draft, Visualization}%
	\projauthor{J.C.-G.}{Juan}{Castillo-Garit}{j.castillo@utem.cl}{0000-0003-0896-9484}%
		{Methodology, Validation, Writing -- review \& editing, Supervision}%
	\projauthor{S.M.-S.}{Samuel}{Montejo-S\'anchez}{smontejo@utem.cl}{0000-0003-1622-3180}%
		{Validation, Writing -- review \& editing, Supervision}%
	\projauthor{R.A.-C.}{Raul}{Arias-Carrasco}{raul.arias@utem.cl}{0000-0003-1698-9282}%
		{Validation, Supervision, Writing -- review \& editing, Visualization}%
}
